\documentclass[11pt,letterpaper]{article}
\usepackage{geometry}
\usepackage{amsmath}
\usepackage{hyperref}
\usepackage{url}
\usepackage{natbib}
\usepackage{booktabs}

\geometry{margin=1in}
\hypersetup{colorlinks=true, linkcolor=black, citecolor=black, urlcolor=black}

\title{Two Recent Methods for Cache Admission: S3-FIFO and Guard}
\author{}
\date{\today}

\begin{document}
\maketitle

\section{S3-FIFO: Quick Demotion via Static FIFO Queues}

S3-FIFO \citep{Yang2023} rethinks cache admission around a single observation: in highly skewed workloads, most objects are accessed only once within any short window, so admitting them into the main cache wastes capacity. The algorithm maintains three static FIFO queues: a small \textit{small} queue $S$ (typically 1\% of cache capacity), a \textit{main} queue $M$, and a \textit{mid} queue $U$. On a cache miss, the new object enters $S$. If it is re-accessed while in $S$, it promotes to $M$. Objects that age out of $S$ without a second access are evicted immediately, achieving what the authors call \textit{quick demotion}.

The admission decision is implicit: an object is admitted to the main cache only after demonstrating at least two accesses. The small queue acts as a probationary filter. S3-FIFO uses no frequency sketches, no learned models, and no adaptive parameters. The queue sizes are fixed fractions of total capacity.

Evaluated on 6594 traces from 14 datasets (856 billion requests), S3-FIFO achieves lower miss ratios than LRU, TinyLFU, and ARC across nearly all traces. On the large cache size, it reduces miss ratios by a mean of 14\% and by more than 32\% on the 10\% worst traces. The simplicity enables 6$\times$ higher throughput than optimized LRU at 16 threads, as FIFO queues avoid the pointer chasing of LRU lists.

\textbf{Strengths.} No parameters to tune. No data structures beyond linked lists. Strong performance on real production traces. Scalable to many cores.

\textbf{Limitations.} No semantic understanding of objects. No uncertainty quantification. Performance degrades on workloads where most objects are accessed exactly twice with the second access falling outside $S$ --- an adversarial pattern the authors acknowledge.

\section{Guard: Robustifying Learning-Augmented Caching}

Guard \citep{Chen2025} addresses a different problem: how to use learned predictions for cache admission without catastrophic failure when the predictor is wrong. Learning-augmented caching algorithms achieve ideal 1-consistency (matching the optimal offline algorithm when predictions are perfect) but can suffer unbounded competitive ratios when predictions are adversarial. Existing robustification methods either sacrifice 1-consistency or introduce $O(k)$ per-request overhead.

Guard wraps any learning-augmented caching algorithm with a lightweight robustification framework. It maintains a confidence score for each prediction and uses a threshold-based strategy: when confidence exceeds a learned threshold, the algorithm follows the prediction; otherwise, it falls back to a safe baseline (e.g., LRU). The key theoretical result is that Guard achieves robustness $2H_{k-1} + 2$ while preserving 1-consistency, with only $O(1)$ additional per-request overhead.

The robustness bound means that even with completely wrong predictions, Guard's cost is at most $2H_{k-1} + 2$ times the optimal offline cost, where $H_{k-1}$ is the $(k-1)$-th harmonic number. This is the best-known trade-off between consistency and robustness for learning-augmented caching.

\textbf{Strengths.} Theoretical guarantees on both consistency and robustness. $O(1)$ overhead preserves the base algorithm's time complexity. Works with any learning-augmented base algorithm.

\textbf{Limitations.} The confidence score requires a trained predictor, which introduces its own assumptions and failure modes. No distribution-free uncertainty guarantees: the confidence score is model-dependent, not statistically calibrated. The robustness bound, while improved, is still a worst-case guarantee that may be loose in practice.

\section{Comparison and Open Directions}

S3-FIFO and Guard represent opposite ends of the design spectrum. S3-FIFO eliminates learned components entirely, achieving strong empirical performance through structural simplicity. Guard embraces learned predictions but wraps them in theoretical safety nets. Neither provides distribution-free uncertainty quantification: S3-FIFO has no uncertainty model, and Guard's confidence scores are model-dependent.

This gap motivates approaches like VUCCA, which use conformal prediction to provide statistically calibrated uncertainty intervals for admission decisions. Conformal prediction offers distribution-free coverage guarantees that neither S3-FIFO's fixed thresholds nor Guard's learned confidence scores provide. The challenge remains to combine the empirical strength of S3-FIFO, the theoretical robustness of Guard, and the statistical guarantees of conformal prediction into a single admission policy.

\bibliographystyle{plainnat}
\bibliography{references}

\end{document}
